\section*{第 \ref{ch:b3:renal-osmoregulation} 章——肾脏生理与渗透调节}

\begin{solution}{exo:b3:renal-osmoregulation:1}
\omterm{def:b3:renal-osmoregulation:nephron}{肾小球}：滤过不含蛋白质的血浆。\omterm{def:b3:renal-osmoregulation:nephron}{近端小管}：重吸收
三分之二的盐分和水、全部葡萄糖与氨基酸、
碳酸氢根。\omterm{def:b3:renal-osmoregulation:nephron}{髓袢}降支：向髓质失去水分。
\omterm{def:b3:renal-osmoregulation:nephron}{髓袢}升支：泵出 NaCl，对水不透过。远端小管：
调节钠（\omterm{def:b3:renal-osmoregulation:controls}{醛固酮}），分泌钾。\omterm{def:b3:renal-osmoregulation:nephron}{集合管}：
在\omterm{prop:b3:renal-osmoregulation:countercurrent}{抗利尿激素}作用下重吸收水，并处理质子与尿素。
\end{solution}

\begin{solution}{exo:b3:renal-osmoregulation:2}
清除率：单位时间内被该物质清除的血浆体积，
即 $UV/P$。此处 $100\times 1.5/2 = \qty{75}{mL/min}$，低于
\omterm{thm:b3:renal-osmoregulation:clearance}{肾小球滤过率} \qty{120}{mL/min}：该物质被滤过，并部分被重吸收
（约为滤过量的 \qty{38}{\%}）。
\end{solution}

\begin{solution}{exo:b3:renal-osmoregulation:3}
淡水鱼体液比环境更咸：水经鳃
通过渗透作用进入体内，因此饮水只会增加水的负担；它
排出大量稀释尿液，其鳃中的\omterm{def:b3:renal-osmoregulation:comparative}{泌氯细胞}从水中
主动摄取 Na$^{+}$ 和 Cl$^{-}$。海水鱼的体液比海水淡：
它经鳃失水，饮入海水以补充，
\omterm{def:b3:renal-osmoregulation:comparative}{泌氯细胞}将吞入的盐分泵出。
\end{solution}

\begin{solution}{exo:b3:renal-osmoregulation:4}
\omterm{prop:b3:renal-osmoregulation:countercurrent}{单效应}是粗升支维持的 \qty{200}{mOsm/L} 差值，
即组织间液与其自身管内液之间的差值，
这一差值靠泵出 NaCl 而不让水随之移动来维持。升支
必须对水不透过，否则水会随盐分移动，差值便会
消失；降支必须对水透过，以使
到达拐弯处的液体已经被组织间液浓缩，
从而每一段都在比前一段更咸的液体上工作。两支
透性相同则无法逐级放大。
\end{solution}

\begin{solution}{exo:b3:renal-osmoregulation:5}
入球小动脉端：$60 - 18 - 28 = \qty{14}{mmHg}$；出球小动脉端：$60 - 18 -
35 = \qty{7}{mmHg}$。滤过水分时，留下的蛋白质
浓缩，胶体\omterm{def:b3:renal-osmoregulation:problem}{渗透压}沿毛细血管升高；若它
达到 \qty{42}{mmHg}，则净压力为零，滤过
在末端之前停止（滤过平衡）。此后\omterm{thm:b3:renal-osmoregulation:clearance}{肾小球滤过率}取决于
血浆流量：流量越快，蛋白质浓缩越慢，因此
滤过在更长的路段上持续进行，\omterm{thm:b3:renal-osmoregulation:clearance}{肾小球滤过率}随
流量而升高。
\end{solution}

\begin{solution}{exo:b3:renal-osmoregulation:6}
滤过负荷：$\qty{40}{mL/min}\times\qty{0.14}{mmol/mL} =
\qty{5.6}{mmol/min}$。排泄量：$70\times 1 = \qty{0.07}{mmol/min}$。
重吸收比例：$5.53/5.6 = \qty{98.75}{\%}$。钠的清除率：$70\times
1/140 = \qty{0.5}{mL/min}$。
\end{solution}

\begin{solution}{exo:b3:renal-osmoregulation:7}
$900/1200 = \qty{0.75}{L}$ 每天。以 \qty{400}{mOsm/L} 计：$900/400 =
\qty{2.25}{L}$ 的尿液，因此约须饮用 $2.25 + 0.9 = \qty{3.15}{L}$
（扣除食物中的水分）：患者的浓缩机制失效，
必须饮水，若无法饮水则有生命危险。
\end{solution}

\begin{solution}{exo:b3:renal-osmoregulation:8}
(a) 尿崩症：尿量 \qtyrange{10}{15}{L}，尿液\omterm{def:b3:renal-osmoregulation:problem}{渗透浓度}
\qtyrange{50}{100}{mOsm/L}，饮水滞后时血浆钠偏高，
口渴不止。(b) \omterm{prop:b3:renal-osmoregulation:countercurrent}{抗利尿激素}分泌不当：尿量少而浓，
水分潴留，血浆钠因稀释而偏低，口渴并未增加，
尽管仍在饮水——意识混乱与惊厥出现于
血浆钠大幅降低之时。(c) \omterm{def:b3:renal-osmoregulation:controls}{醛固酮}过多：钠与水一并潴留，
血容量扩张，血压升高，钾流失（血浆钾
K$^{+}$ 降低），血浆钠接近正常，因水随盐分滞留，肾脏
部分逃逸；尿量正常。(d) 无盐饮食：肾素
与\omterm{def:b3:renal-osmoregulation:controls}{醛固酮}在一天内升高，尿钠降至每天
数毫摩尔，血浆钠维持稳定，血容量与血压略有
降低。
\end{solution}

\begin{solution}{exo:b3:renal-osmoregulation:9}
稀释尿：渗透清除率 $100\times 6/300 = \qty{2}{mL/min}$；
\omterm{prop:b3:renal-osmoregulation:water}{自由水清除率} $6 - 2 = +\qty{4}{mL/min}$：肾脏正在排出
不含溶质的水，如同大量饮水、\omterm{prop:b3:renal-osmoregulation:countercurrent}{抗利尿激素}作用被抑制之后的情形。
浓缩尿：渗透清除率 $1000\times 0.6/300 =
\qty{2}{mL/min}$；\omterm{prop:b3:renal-osmoregulation:water}{自由水清除率} $0.6 - 2 = -\qty{1.4}{mL/min}$：
\qty{1.4}{mL/min} 的自由水正被回送至体内——
脱水，\omterm{prop:b3:renal-osmoregulation:countercurrent}{抗利尿激素}水平高。
\end{solution}

\begin{solution}{exo:b3:renal-osmoregulation:10}
自皮质起给各段编号，设第 $k$ 段的组织间液
为 $I_{k}$。第 $k$ 段的降支液体等于 $I_{k}$；第 $k$ 段的
升支液体则为 $I_{k} - 200$；而第 $k$ 段的升支液体即
是离开拐弯处并依次流经第 $n, n-1, \dots, k+1$ 段的液体。稳态时
在稳态下，以 \qty{300}{} 进入的液体从升支顶端
流出时为 $I_{1} - 200$，且每段组织间液
比上一段高出至多一个\omterm{prop:b3:renal-osmoregulation:countercurrent}{单效应}，故 $I_{k} \le 300 + 200k$，等号在
各段均以满效应工作时成立：拐弯处达到 $300 + 200n$。
因此梯度受限于\omterm{def:b3:renal-osmoregulation:nephron}{髓袢}的长度
（段数）与\omterm{prop:b3:renal-osmoregulation:countercurrent}{单效应}的乘积，而非任何细胞的泵送能力。$n$
是\omterm{def:b3:renal-osmoregulation:nephron}{髓袢}长度除以液体与组织间液达到平衡的距离
所经过的距离。长的近髓\omterm{def:b3:renal-osmoregulation:nephron}{肾单位}\omterm{def:b3:renal-osmoregulation:nephron}{髓袢}，以及
沙漠啮齿动物的极长\omterm{def:b3:renal-osmoregulation:nephron}{髓袢}，给出较大的 $n$；泵的能力
以及直小血管对溶质的冲洗，限制了这一效应。
\end{solution}

\begin{solution}{exo:b3:renal-osmoregulation:11}
\omterm{def:b3:renal-osmoregulation:comparative}{代谢水}：$5\times 0.8 = \qty{4}{g}$ 淀粉，可产生 $4\times 0.6
= \qty{2.4}{g}$。失水：尿液 $3/5500\ \unit{L} = \qty{0.55}{mL}$；
蒸发 \qty{1.5}{g}（已扣除鼻腔的回收）。总
失水 \qty{2.05}{g}，而获得 \qty{2.4}{g}，盈余
\qty{0.35}{g}，足以覆盖少量粪便失水以及种子的
吸湿水是额外收益。它无需饮水即可存活。若无
鼻腔的回收，呼吸失水将为 $1.5/0.3 = \qty{5}{g}$，
它将在几天内死亡。
\end{solution}

\begin{solution}{exo:b3:renal-osmoregulation:12}
在并流中，血液与透析液在膜的中段达到平衡，
驱动扩散的浓度差降至
零：最佳情况下，血液流出时尿素仅剩一半。在逆流中，
最新鲜的透析液与净化最彻底的血液相遇，
梯度沿整个膜面得以维持，血液可
以接近透析液入口的浓度流出——与
\omterm{def:b3:renal-osmoregulation:nephron}{髓袢}的原理相同。清除率：$300\times 0.7 =
\qty{210}{mL/min}$，每周透析 $3\times 4 = \qty{12}{h}$，即
$210\times 720 = \qty{151}{L}$ 每周；两个肾脏以
\qty{125}{mL/min} 清除 $125\times\num{10080} = \qty{1260}{L}$。
透析机约承担肾脏工作量的八分之一，按时间平均
\qty{15}{mL/min}——足以维持生命，却不足以保持健康。
\end{solution}

\begin{solution}{pb:b3:renal-osmoregulation:1}
\textbf{1.} 净压 $55 - 15 - 30 = \qty{10}{mmHg}$；GFR $= 12.5
\times 10 = \qty{125}{mL/min}$。
\textbf{2.} $125\times 1440 = \qty{180}{L/d}$；$180/3 = 60$ 倍。
\textbf{3.} \omterm{thm:b3:renal-osmoregulation:clearance}{菊粉}：$62.5\times 1/0.5 = \qty{125}{mL/min}$，与
GFR 相等。肌酐：$1.3\times 1/0.01 = \qty{130}{mL/min}$，略
大，因为\omterm{def:b3:renal-osmoregulation:nephron}{近端小管}在滤过之外还分泌少量
肌酐。
\textbf{4.} PAH 清除率 $13\times 1/0.02 = \qty{650}{mL/min}$，即
肾血浆流量；血流量 $650/(1 - 0.45) = \qty{1180}{mL/min}$，
约 \qty{1.2}{L/min}；滤过分数 $125/650 = 0.19$。
\textbf{5.} 净压 \qty{15}{mmHg}，GFR \qty{188}{mL/min}（实际
更低，因为胶体\omterm{def:b3:renal-osmoregulation:problem}{渗透压}沿毛细血管升高）。
血管紧张素 II 收缩出球小动脉，升高
\omterm{def:b3:renal-osmoregulation:nephron}{肾小球}压；阻断它则压力与 GFR 都会
略有下降。在糖尿病中，\omterm{def:b3:renal-osmoregulation:nephron}{肾小球}在高压下工作，滤过
膜逐渐受损；降低压力可将\omterm{def:b3:dna-repair:lesions}{损伤}延缓数年。
\textbf{6.} 足细胞及其裂孔膜（或\omterm{def:b3:sensory-systems:cochlea}{基底膜}的
电荷）：截留白蛋白的那一层。每日失去
\qty{3}{g/d} 白蛋白会降低血浆胶体\omterm{def:b3:renal-osmoregulation:problem}{渗透压}，液体从
各处毛细血管\omterm{def:b3:innate-immunity:inflammation}{渗出}，组织肿胀——肾病综合征所致的
水肿。
\textbf{7.} 滤过量 $125\times 5 = \qty{0.625}{mmol/min}$；排泄量
$250\times 1 = \qty{0.25}{mmol/min}$；重吸收 $0.375$，即
\qty{60}{\%}；清除率 $250/5 = \qty{50}{mL/min}$。
\textbf{8.} 滤过量 $180\times 0.14 = \qty{25.2}{mol/d}$ 的 Na$^{+}$，
即 \qty{580}{g} 钠，折合为 NaCl \qty{1.47}{kg}。排泄量 $100\times 1.44
= \qty{144}{mmol/d}$；重吸收率 $1 - 144/\num{25200} = \qty{99.4}{\%}$。
每天将损失近 1.5 千克食盐。
\textbf{9.} 滤过量 $125\times 5 = \qty{0.625}{mmol/min}$，即
\qty{112}{mg/min}，\qty{162}{g/d}；无排泄，因为该值低于
$T_{m}$。滤过负荷达到 $T_{m}$ 时，浓度为 $2.1/0.125 =
\qty{16.8}{mmol/L}$（\qty{300}{mg/dL}）——高于观测到的
阈值约 \qty{10}{mmol/L}，因为各\omterm{def:b3:renal-osmoregulation:nephron}{肾单位}
存在差异，阈值最低者先溢出（滴定曲线的展开）。
\textbf{10.} 滤过量 $125\times 20 = \qty{2.5}{mmol/min}$；排泄量
$2.5 - 2.1 = \qty{0.4}{mmol/min}$，\qty{72}{mg/min}，$576$ mmol，即
每天 \qty{104}{g}。多出的尿量 $576/300 = \qty{1.9}{L}$。葡萄糖
将水分带入尿液（渗透性利尿），患者脱水
并口渴；每天损失的 \qty{104}{g} 葡萄糖相当于 \qty{1.8}{MJ}，
细胞无法摄取葡萄糖，转而燃烧脂肪与蛋白质，致使
体重下降。
\textbf{11.} 重吸收量 $\approx \qty{25}{mol/d}$；ATP $25/3 =
\qty{8.4}{mol/d}$；$8.4\times 50 = \qty{420}{kJ}$，约
占 \qty{7}{MJ} 的 \qty{6}{\%}。
\textbf{12.} 若要分泌每一种废物，小管就需要一种
能识别它的转运蛋白——数以千计，而对从未
遇到过的分子则一个也没有。滤过所有小分子，再重吸收
机体所需的几百种物质，只需为已知有价值的
物质配备转运蛋白；未识别的物质则默认排出。代价是
第~11~题所述的能量。
\textbf{13.} 渗透清除率 $600\times 1/290 = \qty{2.07}{mL/min}$；
\omterm{prop:b3:renal-osmoregulation:water}{自由水清除率} $1 - 2.07 = -\qty{1.07}{mL/min}$：肾脏正在
保留水分。
\textbf{14.} 最小值 $600/1200 = \qty{0.5}{L}$；最大值 $600/50 =
\qty{12}{L}$；按最小值计算的每日需水量为 $0.5 + 0.9 = \qty{1.4}{L}$。
\textbf{15.} \qty{1000}{mOsm} 以 \qty{1200}{mOsm/L} 排出，
需 \qty{0.83}{L} 尿液：未计当日其余溶质时净增 \qty{0.17}{L}。
饮用后：当日溶质总量为 \qty{1600}{mOsm}，需 \qty{1.33}{L}
尿液，加上 \qty{0.9}{L} 不显性失水，共排出 \qty{2.23}{L}，摄入仅 \qty{1}{L}，
亏缺 \qty{1.23}{L}，相比不饮水时的 \qty{1.4}{L}。
一升海水只能使平衡改善 \qty{0.17}{L}。
\textbf{16.} 当日尿液浓度为 \qty{50}{mOsm/L}：$600/50 =
\qty{12}{L}$；每天必须饮用 \qty{12.9}{L}，约 \qty{13}{L}。
\textbf{17.} $B$ 升啤酒带入 $20B$ mOsm 溶质；最大尿量为
$(600 + 20B)/50 = 12 + 0.4B$ 升，且须排出 $B - 0.9$
升：$B - 0.9 \le 12 + 0.4B$，$B \le \qty{21.5}{L}$。约 20
升——但若当日几乎不进食，溶质降至
\qty{200}{mOsm}，饮酒上限降至约 \qty{8}{L}：
饮酒多而不进食者的低钠血症。
\textbf{18.} $298/290 - 1 = \qty{2.8}{\%}$：远超过 \qty{1}{\%}
（引发\omterm{prop:b3:renal-osmoregulation:countercurrent}{抗利尿激素}与口渴的阈值）；两者均已达最大。恒定溶质：
$290\times 42 = 298\times V$，$V = \qty{40.9}{L}$，亏缺约
\qty{1.1}{L}。
\textbf{19.} 6 升水稀释血浆，同时汗液带走
盐分；运动、疼痛和应激会释放\omterm{prop:b3:renal-osmoregulation:countercurrent}{抗利尿激素}，
不论\omterm{def:b3:renal-osmoregulation:problem}{渗透浓度}如何，因而肾脏无法将尿液稀释到足以排出
多余水分。血浆钠下降，脑细胞肿胀——曾引发
惊厥乃至死亡。肾脏需关闭\omterm{prop:b3:renal-osmoregulation:countercurrent}{抗利尿激素}，
并以 $U_{\min}$ 排出 \qty{6}{L} 尿液；跑者应按口渴饮水，
并补充盐分。
\textbf{20.} 每天 $3/5500\ \unit{L} = \qty{0.55}{mL}$。人类肾脏
若排出相同溶质需 $3/1200 = \qty{2.5}{mL}$，而更格卢鼠
只用其五分之一的水。
\textbf{21.} 入：\qty{2.4}{g}。出：$1.6 + 0.3 + 0.55 = \qty{2.45}{g}$。
每天平衡误差在 \qty{0.05}{g} 之内，这由种子的
吸湿水（约占 \qty{5}{g} 的百分之几）补足。
\textbf{22.} 水量 $500\times 0.65 = \qty{325}{L}$；四分之一为
\qty{81}{L}；按 \qty{5}{L/d}，可用 16 天。
\textbf{23.} $150/800 = \qty{0.19}{L}$，约 \qty{190}{mL} 盐水
——少于 \qty{100}{mL} 海水与鱼体的水分
之和，因此这只鸟获得了水分。其肾脏在 \qty{600}{mOsm/L}
（NaCl \qty{300}{mmol/L}），排出相同盐分需 \qty{0.5}{L} 尿液，
所排水量多于摄入量：肾脏无法超过海水的浓度，因而鸟不能饮用海水，
而\omterm{def:b3:renal-osmoregulation:comparative}{盐腺}的浓度为海水的两倍，
则可以。
\textbf{24.} 摄入水 \qty{30}{g}，故每天尿量约 \qty{30}{mL}，
\omterm{def:b3:renal-osmoregulation:problem}{渗透浓度}仅为血液的很小一部分（几十
mOsm/L，对比 \num{300}）。盐分经该尿液排出，
并经鳃扩散；若无\omterm{def:b3:renal-osmoregulation:comparative}{泌氯细胞}从
含每升 1 毫摩尔盐的水中主动摄取，鱼体将
在数日内耗竭。
\textbf{25.} GFR \qty{125}{mL/min}，每天 \qty{180}{L}；必需
尿量每天 \qty{0.5}{L}；一升海水
排出盐分后净得 \qty{0.17}{L}——几乎等于零。
\end{solution}
